\section*{Capítulo \ref{ch:b1:linmaps} --- Aplicações Lineares}

\begin{solution}{exo:b1:linmaps:1}
(1) \omterm{def:b1:linmaps:def}{Linear}: as \omterm{prop:b1:vspaces:coordinates}{coordenadas} são expressões \omterm{def:b1:linmaps:def}{lineares}. (2) Não \omterm{def:b1:linmaps:def}{linear}:
$u(2(1,1)) = 4 \neq 2 = 2u(1,1)$. (3) \omterm{def:b1:linmaps:def}{Linear}: derivar e multiplicar por $X$ o são, e somas de
\omterm{def:b1:linmaps:def}{aplicações lineares} também. (4) \omterm{def:b1:linmaps:def}{Linear}: a avaliação respeita as
operações ponto a ponto.
\end{solution}

\begin{solution}{exo:b1:linmaps:2}
\omterm{def:b1:linmaps:kerim}{Núcleo}: resolva $x + y - z = 0$, $2x + y + z = 0$, $3x + 2y = 0$. Da terceira, $y = -\frac{3x}{2}$; a primeira
dá $z = x + y =
-\frac x2$; verifique na segunda: $2x - \frac{3x}{2} - \frac x2 = 0$: satisfeita. Logo $\ker u = \operatorname{Vect}\bigl((2, -3, -1)\bigr)$
(tomando $x = 2$), dimensão $1$.

Pelo teorema do \omterm{def:b1:linmaps:kerim}{núcleo} e da \omterm{def:b1:linmaps:kerim}{imagem}: $\operatorname{rk} u = 3 - 1 = 2$. \omterm{def:b1:linmaps:kerim}{Imagem}: gerada pelas
imagens da \omterm{def:b1:vspaces:free}{base} canônica, $u(e_1) = (1,2,3)$, $u(e_2) =
(1,1,2)$, $u(e_3) = (-1,1,0)$; as duas primeiras são
\omterm{def:b1:vspaces:free}{livres}, e o \omterm{def:b1:linmaps:rank}{posto} é $2$: \omterm{def:b1:vspaces:free}{base} $\bigl((1,2,3), (1,1,2)\bigr)$.

Não \omterm{def:b1:logic:inj}{injetiva} ($\ker \neq \{0\}$), não \omterm{def:b1:logic:inj}{sobrejetiva} (\omterm{def:b1:linmaps:rank}{posto} $2 < 3$): coerente com
o \cref{cor:b1:linmaps:samedim}.
\end{solution}

\begin{solution}{exo:b1:linmaps:3}
\emph{\omterm{def:b1:linmaps:kerim}{Núcleo}:} $P = P'$ força $\deg P = \deg P'$ a não ser que $P = 0$; mas $\deg P' < \deg P$
para $P \neq 0$: logo $\ker u = \{0\}$, e $u$, endomorfismo \omterm{def:b1:logic:inj}{injetivo} do espaço de
\omterm{def:b1:findim:def}{dimensão finita} $\R_n[X]$, é um isomorfismo (o \cref{cor:b1:linmaps:samedim}).

\emph{Inversa:} seja $v(P) = P + P' + P'' + \dots + P^{(n)}$ (uma soma finita em $\R_n[X]$). Então
\[
v\bigl(u(P)\bigr) = \sum_{k=0}^{n} (P - P')^{(k)}
= \sum_{k=0}^{n} P^{(k)} - \sum_{k=0}^{n} P^{(k+1)}
= P - P^{(n+1)} = P ,
\]
telescopando, pois $P^{(n+1)} = 0$. Logo $v = u^{-1}$.
\end{solution}

\begin{solution}{exo:b1:linmaps:4}
$\operatorname{im}(v \circ u) = v(\operatorname{im} u) \subseteq
\operatorname{im} v$: \omterm{def:b1:linmaps:rank}{posto} $\leq \operatorname{rk} v$. E $v$ restrita a $\operatorname{im} u$ tem \omterm{def:b1:linmaps:kerim}{imagem} $\operatorname{im}(vu)$, com o
teorema do \omterm{def:b1:linmaps:kerim}{núcleo} e da \omterm{def:b1:linmaps:kerim}{imagem} dentro de $\operatorname{im} u$: $\operatorname{rk}(vu) \leq
\dim\operatorname{im} u = \operatorname{rk} u$.
\end{solution}

\begin{solution}{exo:b1:linmaps:5}
$u^2 = 0$ significa $u(u(x)) = 0$ para todo $x$: todo $u(x)$ está em
$\ker u$, isto é, $\operatorname{im} u \subseteq \ker u$. Então, pelo teorema do \omterm{def:b1:linmaps:kerim}{núcleo} e da \omterm{def:b1:linmaps:kerim}{imagem}:
\[
\dim E = \dim\ker u + \operatorname{rk} u \geq 2\operatorname{rk} u .
\]
Exemplo de igualdade em $\R^2$: $u(x, y) = (y, 0)$: $u^2 = 0$, $\operatorname{rk} u = 1 = \frac{\dim E}{2}$.
\end{solution}

\begin{solution}{exo:b1:linmaps:6}
$(pq)^2 = pqpq = ppqq = pq$ (comutação): uma \omterm{def:b1:linmaps:projection}{projeção} (o \cref{thm:b1:linmaps:projchar}).

\omterm{def:b1:linmaps:kerim}{Imagem}: $\operatorname{im}(pq) \subseteq \operatorname{im} p$, e $= \operatorname{im}(qp) \subseteq \operatorname{im} q$: contida na interseção. Reciprocamente, se
$x \in \operatorname{im} p \cap
\operatorname{im} q$, então $p(x) = x$ e $q(x) = x$ (os pontos fixos caracterizam a \omterm{def:b1:linmaps:kerim}{imagem} de
uma \omterm{def:b1:linmaps:projection}{projeção}), logo $pq(x) = x$: $x \in
\operatorname{im}(pq)$.

\omterm{def:b1:linmaps:kerim}{Núcleo}: $\ker p \subseteq \ker(qp) = \ker(pq)$ e do mesmo modo $\ker q
\subseteq \ker(pq)$: a soma está contida.
Reciprocamente, seja $pq(x) =
0$, e escreva
\[
x = \underbrace{q(x)}_{\in\, \ker p} +
\underbrace{(x - q(x))}_{\in\, \ker q} :
\]
o primeiro termo satisfaz $p(q(x)) = 0$, logo está em $\ker p$; o segundo
está em $\ker q$ pois $q(x - q(x)) = q(x) - q^2(x) = 0$. Logo $x \in \ker p + \ker q$.
\end{solution}

\begin{solution}{exo:b1:linmaps:7}
Note primeiro as inclusões gerais $\ker u \subseteq \ker u^2$ e $\operatorname{im} u^2 \subseteq \operatorname{im} u$, e, pelo teorema do
\omterm{def:b1:linmaps:kerim}{núcleo} e da \omterm{def:b1:linmaps:kerim}{imagem}, (2) $\iff$ (3) (\omterm{def:b1:linmaps:kerim}{núcleos} iguais $\iff$ \omterm{def:b1:linmaps:rank}{postos}
iguais $\iff$ imagens iguais, dadas as inclusões).

(1 $\Rightarrow$ 2): seja $u^2(x) = 0$; então $u(x) \in \ker u \cap
\operatorname{im} u = \{0\}$, logo $x \in \ker u$.

(2 $\Rightarrow$ 1): por Grassmann e pelo teorema do \omterm{def:b1:linmaps:kerim}{núcleo} e da \omterm{def:b1:linmaps:kerim}{imagem},
$\dim(\ker u +
\operatorname{im} u) = \dim\ker u + \operatorname{rk} u -
\dim(\ker u \cap \operatorname{im} u) = \dim E - \dim(\ker u \cap
\operatorname{im} u)$: a soma é $E$ se e somente se a interseção é $\{0\}$. Seja
$y \in \ker u \cap \operatorname{im} u$: $y = u(x)$ e $u(y) = 0$, logo $u^2(x) = 0$, e então (por (2)) $u(x) = 0$: $y = 0$.
Logo $E = \ker u \oplus \operatorname{im} u$.
\end{solution}

\begin{solution}{exo:b1:linmaps:8}
Se $\varphi = 0$: então $\ker\varphi = E \subseteq \ker\psi$ força $\psi = 0 = 0\cdot\varphi$. Se $\varphi \neq 0$: $\ker\varphi$ é um \omterm{def:b1:linmaps:forms}{hiperplano};
escolha $a \notin \ker\varphi$. Ponha $\lambda =
\frac{\psi(a)}{\varphi(a)}$. A forma $\psi - \lambda\varphi$ se anula em $\ker\varphi$ (ambas
se anulam, pela inclusão) e em $a$: ela se anula em $\ker\varphi \oplus Ka = E$. Logo
$\psi =
\lambda\varphi$.
\end{solution}

\begin{solution}{exo:b1:linmaps:9}
Suponha $\lambda_0 x + \lambda_1 u(x) + \dots + \lambda_{n-1}
u^{n-1}(x) = 0$. Aplique $u^{n-1}$: todos os termos com um fator
$u^{\geq
n}$ morrem, restando $\lambda_0 u^{n-1}(x) = 0$, logo $\lambda_0 = 0$. Aplique $u^{n-2}$ à
relação restante: $\lambda_1 u^{n-1}(x) =
0$, logo $\lambda_1 = 0$; e assim por diante. A família é
\omterm{def:b1:vspaces:free}{livre}; tendo tamanho $n = \dim E$, é uma \omterm{def:b1:vspaces:free}{base} (a \cref{prop:b1:findim:twoofthree}). (Nesta \omterm{def:b1:vspaces:free}{base}, $u$
age como um deslocamento --- o modelo da nilpotência maximal.)
\end{solution}

\begin{solution}{exo:b1:linmaps:10}
\begin{enumerate}
  \item $f \circ f = -\mathrm{id}$ é \omterm{def:b1:logic:inj}{bijetiva}, logo $f$ também é (a \cref{prop:b1:logic:comp} adaptada: $f$
        tem a inversa bilateral $-f$). Se $f(x) = \lambda x$ com $x \neq 0$:
        aplicando $f$, $-x = \lambda^2 x$, logo $\lambda^2 = -1$: impossível em $\R$.
  \item Construa a família de modo guloso. Tome $x_1 \neq 0$: $(x_1,
        f(x_1))$ é
        \omterm{def:b1:vspaces:free}{livre} por (1). Se $\operatorname{Vect}$ da família corrente $\bigl(x_1, f(x_1), \dots, x_k,
        f(x_k)\bigr)$, que
        chamamos $V_k$ --- um \omterm{def:b1:vspaces:subspace}{subespaço} estável por $f$ (cada
        gerador é levado noutro gerador ou no seu oposto: $f(f(x_i)) = -x_i$)
        --- não é todo o $E$, escolha $x_{k+1}
        \notin V_k$. \emph{Afirmação: a
        família aumentada é \omterm{def:b1:vspaces:free}{livre}.} Suponha $\alpha x_{k+1} + \beta f(x_{k+1}) + v = 0$ com $v
        \in V_k$ e
        $(\alpha, \beta) \neq (0,0)$. Aplique $f$: $\alpha f(x_{k+1}) - \beta x_{k+1} + f(v) = 0$ com $f(v)
        \in V_k$. Elimine $f(x_{k+1})$
        entre as duas relações (multiplique a primeira por
        $\alpha$, a segunda por $-\beta$, e some):
        \[
        (\alpha^2 + \beta^2)\, x_{k+1} \in V_k ,
        \]
        e $\alpha^2 + \beta^2 \neq 0$ força $x_{k+1} \in V_k$: contradição. Logo a construção
        continua, acrescentando vetores \emph{dois a dois}, até
        $V_k = E$: a família final é uma \omterm{def:b1:vspaces:free}{base} de tamanho par, e $n$
        é par.
\end{enumerate}
\end{solution}

\begin{solution}{exo:b1:linmaps:11}
\omterm{def:b1:reals:bounds}{Cota superior}: para todo $x$, $(u + v)(x) = u(x) + v(x) \in
\operatorname{im} u + \operatorname{im} v$, logo
\[
\operatorname{rk}(u + v)
\leq \dim(\operatorname{im} u + \operatorname{im} v)
\leq \operatorname{rk} u + \operatorname{rk} v
\]
(Grassmann, o \cref{thm:b1:findim:grassmann}). Cota inferior: aplique a \omterm{def:b1:reals:bounds}{cota superior} ao par
$(u + v, -v)$, cuja soma é $u$:
\[
\operatorname{rk} u \leq \operatorname{rk}(u + v) +
\operatorname{rk}(-v) = \operatorname{rk}(u + v) +
\operatorname{rk} v,
\]
logo $\operatorname{rk} u - \operatorname{rk} v \leq
\operatorname{rk}(u+v)$; trocar $u$ e $v$ dá o valor absoluto.
\end{solution}

\begin{solution}{exo:b1:linmaps:12}
\emph{Fórmula exata.} Seja $v'$ a restrição de $v$ ao \omterm{def:b1:vspaces:subspace}{subespaço}
$\operatorname{im} w$. A sua \omterm{def:b1:linmaps:kerim}{imagem} é $v(w(F)) =
\operatorname{im}(v \circ w)$, e o seu \omterm{def:b1:linmaps:kerim}{núcleo} é $\ker v \cap
\operatorname{im} w$. O teorema do
\omterm{def:b1:linmaps:kerim}{núcleo} e da \omterm{def:b1:linmaps:kerim}{imagem} para $v'$ no espaço $\operatorname{im} w$:
\[
\operatorname{rk} w = \dim\operatorname{im} w
= \operatorname{rk}(v \circ w) + \dim(\ker v \cap
\operatorname{im} w) .
\]

\emph{Frobenius.} Aplique a fórmula exata duas vezes, a $w$ e a
$w \circ u$:
\[
\operatorname{rk} w - \operatorname{rk}(vw)
= \dim\bigl(\ker v \cap \operatorname{im} w\bigr),
\qquad
\operatorname{rk}(wu) - \operatorname{rk}(vwu)
= \dim\bigl(\ker v \cap \operatorname{im}(wu)\bigr) .
\]
Como $\operatorname{im}(w \circ u) \subseteq \operatorname{im}
w$, a segunda interseção está contida na primeira, e a sua
dimensão não é maior:
\[
\operatorname{rk}(wu) - \operatorname{rk}(vwu)
\;\leq\; \operatorname{rk} w - \operatorname{rk}(vw) ,
\]
o que se rearranja na desigualdade de Frobenius. Com $w =
\mathrm{id}_F$ (\omterm{def:b1:linmaps:rank}{posto}
$\dim F$, e $\operatorname{im}\,
\mathrm{id}_F = F$): $\operatorname{rk} v + \operatorname{rk} u
\leq \dim F + \operatorname{rk}(vu)$, a desigualdade de Sylvester ---
demonstrada de novo, em forma matricial, no \cref{exo:b1:matrices:10}.
\end{solution}

\begin{solution}{pb:b1:linmaps:1}
\textbf{1.} $(\mathrm{id} - p)^2 = \mathrm{id} - 2p + p^2 =
\mathrm{id} - p$: um projetor. Se $y = x - p(x)$, então $p(y) =
p(x) - p^2(x) = 0$, e
reciprocamente $x \in \ker p$ dá $x =
(\mathrm{id} - p)(x)$: $\operatorname{im}(\mathrm{id} - p) = \ker
p$. E $(\mathrm{id} - p)(x) = 0 \iff p(x) = x \iff x \in
\operatorname{im} p$ (pontos fixos,
o \cref{thm:b1:linmaps:projchar}): $\ker(\mathrm{id} - p) =
\operatorname{im} p$.

\textbf{2.} $(\lambda\,\mathrm{id} + \mu p)^2 =
\lambda^2\,\mathrm{id} + (2\lambda\mu + \mu^2)\,p$. O par $(\mathrm{id}, p)$ é \omterm{def:b1:vspaces:free}{livre} em $\mathcal{L}(E)$: $p = c\,
\mathrm{id}$ daria
$c^2 = c$, logo $p = 0$ ou $\mathrm{id}$, o que está excluído.
Identificando coeficientes, a \omterm{def:b1:logic:map}{aplicação} é um projetor se e somente
se $\lambda^2 = \lambda$ e $2\lambda\mu + \mu^2 = \mu$. Para $\lambda = 0$: $\mu \in \{0, 1\}$. Para $\lambda = 1$: $\mu^2 +
\mu = 0$, $\mu \in \{0, -1\}$.
Exatamente quatro projetores no plano: $0$, $p$, $\mathrm{id}$, $\mathrm{id} - p$.

\textbf{3.} $(\lambda\,\mathrm{id} + \mu p)(\lambda'\,\mathrm{id}
+ \mu' p) = \lambda\lambda'\,\mathrm{id} + (\lambda\mu' +
\mu\lambda' + \mu\mu')\,p$: o plano é estável por composição. Como $p^k = p$
para todo $k \geq 1$, para $Q =
\sum_k a_k X^k$:
\[
Q(p) = a_0\,\mathrm{id} + \Bigl(\sum_{k \geq 1} a_k\Bigr) p
= Q(0)\,\mathrm{id} + \bigl(Q(1) - Q(0)\bigr)\,p .
\]

\textbf{4.} Em $\operatorname{im} p$ (onde $p$ age como a identidade), $\lambda\,\mathrm{id} + \mu p$
multiplica por $\lambda
+ \mu$; em $\ker p$, por $\lambda$. Como $E = \operatorname{im} p
\oplus \ker p$, a
\omterm{def:b1:logic:map}{aplicação} é \omterm{def:b1:logic:inj}{bijetiva} se e somente se $\lambda \neq 0$ e $\lambda + \mu \neq 0$. Resolvendo
$\lambda\alpha = 1$, $\lambda\beta + \mu\alpha + \mu\beta = 0$ na regra de composição da questão 3:
\[
(\lambda\,\mathrm{id} + \mu p)^{-1}
= \frac1\lambda\,\mathrm{id} -
\frac{\mu}{\lambda(\lambda + \mu)}\,p ,
\]
cuja ação é por $1/\lambda$ em $\ker p$ e por $1/(\lambda +
\mu)$ em $\operatorname{im} p$, como tem de
ser.

\textbf{5.} Escreva $F = \operatorname{im} p = \operatorname{im}
p'$. Para todo $x$, $p'(x) \in F$ e $p$ fixa $F$
ponto a ponto: $p(p'(x)) = p'(x)$, isto é, $p\,p' = p'$; simetricamente $p'\,p =
p$. Quando
duas projeções partilham a \omterm{def:b1:linmaps:kerim}{imagem}, aquela aplicada \emph{primeiro}
decide: a sua saída já está em $F$, onde a \omterm{def:b1:linmaps:projection}{projeção} externa age
como a identidade e nada muda.

\textbf{6.} $(p + q)^2 = p^2 + pq + qp + q^2 = (p + q) + pq +
qp$, logo, sendo $p + q$ um projetor, $pq + qp = 0$. Componha
à esquerda com $p$: $pq + pqp = 0$; à direita com $p$: $pqp + qp
= 0$. Subtraindo,
$pq = qp$; então $pq + qp = 2pq = 0$ e a característica não é $2$: $pq = qp = 0$.

\textbf{7.} Com $pq = qp = 0$, o mesmo desenvolvimento dá $(p +
q)^2 = p + q$. \omterm{def:b1:linmaps:kerim}{Imagem}:
$\operatorname{im}(p + q) \subseteq
\operatorname{im} p + \operatorname{im} q$ sempre. Reciprocamente, para $x \in \operatorname{im} p$: $q(x) = q(p(x)) = 0$, logo $(p +
q)(x) = p(x) = x$ e $x \in \operatorname{im}(p+q)$;
o mesmo para $\operatorname{im} q$. Caráter direto: $x \in \operatorname{im} p
\cap \operatorname{im} q$ dá $x = p(x) = p(q(x)) = 0$. \omterm{def:b1:linmaps:kerim}{Núcleo}: se
$p(x) + q(x) = 0$, aplicar $p$ dá $p(x) +
p(q(x)) = p(x) = 0$, e aplicar $q$ dá $q(x) = 0$: $\ker(p
+ q) = \ker p \cap \ker q$ (a
inclusão recíproca é clara).

\textbf{8.} $p - q$ é um projetor se e somente se $\mathrm{id} - (p - q) =
(\mathrm{id} - p) + q$ o for
(questão 1 duas vezes). Pelas questões 6--7 aplicadas aos
projetores $\mathrm{id} - p$ e $q$, isto vale se e somente se $(\mathrm{id} - p)q = q(\mathrm{id} - p) = 0$, isto é, se
e somente se $pq = q$ e $qp = q$.

\textbf{9.} $pq = q$ significa que $p$ fixa todo $q(x)$, isto é,
$\operatorname{im} q \subseteq \ker(p - \mathrm{id}) =
\operatorname{im} p$. E $qp = q$ significa $q\bigl((\mathrm{id} -
p)(x)\bigr) = 0$ para todo $x$, isto é, $q$ se
anula em $\operatorname{im}(\mathrm{id} - p) = \ker p$: $\ker p \subseteq
\ker q$. Os dois passos são equivalências: a ordem
$q \leq p$ diz que $q$ projeta sobre uma \omterm{def:b1:linmaps:kerim}{imagem} menor, ao longo de um
\omterm{def:b1:linmaps:kerim}{núcleo} maior.

\textbf{10.} Desenvolvendo, $(\mathrm{id} - p)(\mathrm{id} - q) =
\mathrm{id} - p - q + pq = \mathrm{id} - r$. Os projetores $\mathrm{id} - p$ e $\mathrm{id} - q$
comutam, logo, pelo \cref{exo:b1:linmaps:6}, o seu produto $\mathrm{id} - r$ é o projetor sobre
$\operatorname{im}(\mathrm{id} - p) \cap
\operatorname{im}(\mathrm{id} - q) = \ker p \cap \ker q$ ao longo de $\ker(\mathrm{id} - p) + \ker(\mathrm{id} - q) =
\operatorname{im} p + \operatorname{im} q$. Pela questão 1, $r =
\mathrm{id} - (\mathrm{id} - r)$ é então o projetor
com $\operatorname{im} r = \operatorname{im} p + \operatorname{im}
q$ e $\ker r = \ker p \cap \ker q$.

\textbf{11.} $p_i$ está bem definido (unicidade da decomposição) e
é \omterm{def:b1:linmaps:def}{linear} (a decomposição de $x + \lambda y$ é a soma das decomposições, de
novo por unicidade). Para $x_i
\in F_i$ a decomposição é o próprio $x_i$,
logo $p_i(x_i) = x_i$: $p_i^2 = p_i$, e $p_j(x_i) = 0$ para $j \neq i$: $p_i p_j = 0$ ($p_j(x) \in F_j$). Somando as
componentes, $\sum_i p_i =
\mathrm{id}$. Por fim, $\operatorname{im} p_i = F_i$ e $\ker
p_i = \bigoplus_{j \neq i} F_j$.

\textbf{12.} $p_i = p_i \circ \mathrm{id} = p_i\sum_j p_j =
p_i^2 + \sum_{j \neq i} p_i p_j = p_i^2$: cada $p_i$ é um projetor. Todo $x = \mathrm{id}(x) = \sum_i p_i(x)$ está em
$\sum_i \operatorname{im} p_i$: as imagens somam $E$. Caráter direto: suponha $y_1 + \dots + y_k = 0$ com
$y_i \in
\operatorname{im} p_i$, de modo que $p_i(y_i) = y_i$. Aplique $p_j$: $p_j
(y_i) = p_j p_i (y_i) = 0$ para $i \neq j$, logo
$0 = p_j\bigl(\sum
y_i\bigr) = y_j$, para todo $j$. Portanto $E = \bigoplus_i
\operatorname{im} p_i$.

\textbf{13.} $q = \mathrm{id} - p$, e a questão 1 dá $pq =
p - p^2 = 0 = qp$ diretamente: para dois
projetores, somar a identidade já força a ortogonalidade do par.

\textbf{14.} $p + q = \mathrm{id} - r$ com $r$ um projetor, e $(\mathrm{id} - r)$ é um projetor
(questão 1): logo $p + q$ é um projetor, e a questão 6 dá $pq = qp = 0$.
Por simetria ($q +
r = \mathrm{id} - p$ e $p + r = \mathrm{id} - q$), todos os produtos dois a dois se
anulam, e a questão 12 conclui: $E =
\operatorname{im} p \oplus \operatorname{im} q \oplus
\operatorname{im} r$, automaticamente.

\textbf{15.} \emph{Lema.} Por indução com Grassmann (o \cref{thm:b1:findim:grassmann}):
\[
\dim(F_1 + \dots + F_k) \leq \dim(F_1 + \dots + F_{k-1}) + \dim
F_k \leq \dots \leq \sum_i \dim F_i .
\]
Se o total é uma igualdade, cada passo o é: $(F_1 + \dots +
F_{j-1}) \cap F_j = \{0\}$ para todo $j$, e
uma relação $y_1 +
\dots + y_k = 0$ ($y_i \in F_i$) colapsa a partir da direita: $y_k
\in (F_1 + \dots + F_{k-1}) \cap F_k = \{0\}$,
depois $y_{k-1} =
0$, etc.: a soma é direta. Reciprocamente, uma \omterm{def:b1:vspaces:sum}{soma
direta} tem dimensões aditivas (concatene \omterm{def:b1:vspaces:free}{bases}). \emph{Aplicação:}
$x = \sum_i p_i(x)$ mostra que $E = \sum_i \operatorname{im} p_i$, logo $n \leq \sum_i \operatorname{rk} p_i$; a hipótese dá a igualdade, e
portanto o caráter direto. Produtos: fixe $j$ e $y \in
\operatorname{im} p_j$. Então $y = \sum_i p_i(y)$
com $p_i(y)
\in \operatorname{im} p_i$, ao passo que $y = y$ é também uma decomposição (só a
componente $j$); a unicidade força $p_i(y)
= 0$ para $i \neq j$. Aplicado a
$y = p_j(x)$: $p_i p_j = 0$.

\textbf{16.} Se $u^k(x) = 0$ então $u^{k+1}(x) = u(0) = 0$. E $\operatorname{im} u^{k+1} = u^k\bigl(u(E)\bigr) \subseteq
u^k(E) = \operatorname{im} u^k$.

\textbf{17.} Suponha $\ker u^{r} = \ker u^{r+1}$ e seja $x \in
\ker u^{r+2}$: então $u(x) \in \ker u^{r+1} = \ker u^{r}$, logo $u^{r+1}(x) = 0$:
$x \in \ker u^{r+1}$. Com a questão 16, $\ker u^{r+1} = \ker u^{r+2}$, e por indução todos os \omterm{def:b1:linmaps:kerim}{núcleos}
posteriores coincidem com $\ker u^{r}$. Para as imagens: o teorema do
\omterm{def:b1:linmaps:kerim}{núcleo} e da \omterm{def:b1:linmaps:kerim}{imagem} dá $\dim\operatorname{im} u^k = n - \dim\ker u^k$, logo as dimensões das imagens congelam
exatamente quando as dos \omterm{def:b1:linmaps:kerim}{núcleos} congelam, e, com as inclusões da
questão 16, dimensões iguais significam \omterm{def:b1:vspaces:subspace}{subespaços} iguais
(o \cref{thm:b1:findim:subspaces}).

\textbf{18.} A sequência $\bigl(\dim\ker u^k\bigr)_k$ é não decrescente com valores em
$\intint{0}{n}$; ela não pode crescer estritamente $n + 1$ vezes, logo algum
$\dim\ker u^{r} = \dim\ker
u^{r+1}$ com $r \leq n$, donde $\ker u^{r} = \ker u^{r+1}$ (inclusão mais igualdade de
dimensões). Tome $r$ mínimo; a questão 17 congela tudo a partir de
$r$, imagens incluídas.

\textbf{19.} Interseção: seja $x \in \ker u^{r} \cap
\operatorname{im} u^{r}$, digamos $x = u^{r}(y)$ com $u^{r}(x) =
0$. Então
$u^{2r}(y) = 0$, e $\ker u^{2r} = \ker u^{r}$ (questão 17), logo $x = u^{r}(y) = 0$. Dimensões: o teorema do
\omterm{def:b1:linmaps:kerim}{núcleo} e da \omterm{def:b1:linmaps:kerim}{imagem} para $u^{r}$ dá $\dim\ker u^{r} + \dim\operatorname{im} u^{r}
= n$; com interseção trivial,
Grassmann faz da soma um \omterm{def:b1:vspaces:subspace}{subespaço} de dimensão $n$: $E = \ker u^{r} \oplus
\operatorname{im} u^{r}$.

\textbf{20.} Estabilidade: $u^{r}(u(x)) = u(u^{r}(x)) = 0$ para $x
\in \ker u^{r}$; e $u(u^{r}(y)) = u^{r}(u(y)) \in
\operatorname{im} u^{r}$. Em $N = \ker u^{r}$:
$(u|_N)^{r} = 0$ por definição de $N$: nilpotente. Em $I = \operatorname{im}
u^{r}$: $\ker(u|_I) = \ker u \cap I \subseteq \ker u^{r} \cap I
= \{0\}$, logo
$u|_I$ é um endomorfismo \omterm{def:b1:logic:inj}{injetivo} do espaço de \omterm{def:b1:findim:def}{dimensão finita} $I$,
e portanto \omterm{def:b1:logic:inj}{bijetivo} (o \cref{cor:b1:linmaps:samedim}).

\textbf{21.} Seja $x = a + b$ com $a \in N$, $b \in I$. Então $u(x) = u(a) + u(b)$ com
$u(a) \in N$ e $u(b) \in I$ (questão 20): esta \emph{é} a decomposição de $u(x)$,
logo $\pi(u(x)) = u(a) = u(\pi(x))$: $\pi u = u\pi$.

\textbf{22.} $u^2(x,y,z) = u(y, 0, z) = (0, 0, z)$ e $u^3(x,y,z) = u(0,0,z) = (0,0,z) = u^2(x,y,z)$. \omterm{def:b1:linmaps:kerim}{Núcleos}: $\ker u
= \{y = z = 0\} = \operatorname{Vect}(e_1)$, $\ker u^2 = \{z =
0\} = \operatorname{Vect}(e_1, e_2)$, $\ker u^3 = \ker u^2$:
estabilização em $r = 2$. Imagens: $\operatorname{im} u =
\operatorname{Vect}(e_1, e_3)$, $\operatorname{im} u^2 =
\operatorname{Vect}(e_3)$. Fitting: $\R^3 =
\operatorname{Vect}(e_1, e_2) \oplus \operatorname{Vect}(e_3)$,
e $\pi(x, y, z) = (x, y, 0)$. Verificação: $\pi u(x,y,z) = \pi(y, 0,
z) = (y, 0, 0)$ e $u\pi(x,y,z) = u(x, y, 0) = (y, 0, 0)$: iguais. No primeiro fator
$u(x, y, 0) = (y, 0, 0)$, cujo quadrado é $0$: nilpotente; no segundo $u(0,0,z) = (0,0,z)$: a
identidade, \omterm{def:b1:logic:inj}{bijetiva}.

\textbf{23.} Se $u^m = 0$ então $\ker u^m = E$; como os \omterm{def:b1:linmaps:kerim}{núcleos} estão
congelados a partir de $r$, $\ker u^{r} = \ker u^{\max(m, r)} = E$. Reciprocamente, $\ker u^{r} = E$ significa
$u^{r} = 0$. E $\ker u^{r} =
E \iff$ o projetor de Fitting é sobre $E$ ao longo de
$\{0\}$, isto é, $\pi = \mathrm{id}$. Por fim, $r \leq n$ (questão 18) dá: todo
endomorfismo nilpotente satisfaz $u^{n} = 0$ --- o índice de nilpotência
nunca ultrapassa a dimensão.

\textbf{24.} Seja $m$ um índice de nilpotência de $u|_A$: $A
\subseteq \ker u^{m} \subseteq \ker u^{\max(m,r)} = \ker u^{r}$.
Como $u|_B$ é \omterm{def:b1:logic:inj}{bijetiva}, $B = u(B) = u^{k}(B) \subseteq
\operatorname{im} u^{k}$ para todo $k$, em particular $B
\subseteq \operatorname{im} u^{r}$.
Então
\[
n = \dim A + \dim B \leq \dim\ker u^{r} +
\dim\operatorname{im} u^{r} = n :
\]
as duas inclusões são igualdades de dimensões, e portanto de
\omterm{def:b1:vspaces:subspace}{subespaços}: $A = \ker u^{r}$, $B = \operatorname{im} u^{r}$.

\textbf{25.} (i) A Parte III é um dicionário: as decomposições
$E = F_1
\oplus \dots \oplus F_k$ correspondem exatamente a famílias de projetores com $\sum p_i = \mathrm{id}$ e
$p_i p_j = 0$, sendo os $F_i$ as imagens. (ii) Para $k = 3$ os
complementares $\mathrm{id} - p_i$ são eles próprios projetores, o que fechou o
argumento sem hipótese adicional; para $k$ geral é preciso $\sum_i \operatorname{rk} p_i \leq n$,
uma desigualdade que o traço do segundo ano dá de graça ($\operatorname{tr} p = \operatorname{rk}
p$ para
um projetor, e os traços somam $\operatorname{tr}
\mathrm{id} = n$). (iii) O \cref{exo:b1:linmaps:7} é o lema de
Fitting no caso já estabilizado $r \leq 1$; em geral deixa-se as cadeias
de \omterm{def:b1:linmaps:kerim}{núcleos} e de imagens congelarem, o que leva no máximo $n$
passos. (iv) No volume do segundo ano de graduação, aplicado a
$u - \lambda\,
\mathrm{id}$, o fator nilpotente torna-se o autoespaço generalizado em
$\lambda$ e os projetores da Parte III tornam-se os projetores
espectrais da teoria da redução. O teorema da Parte IV é o
\emph{lema de Fitting}.
\end{solution}
